% Generated by ops/build_paper_hosted_level_averages.py; do not edit.
\paragraph{Evaluation parity across deployments.}
One column of this cohort rests on wording the other deployments did not
answer: Claude Opus 5's Python cells in Fig.~\ref{fig:hosted-verified-breadth}
and Table~\ref{tab:hosted-level-averages-medium} use the revised condition of
App.~\ref{app:hosted-python-prompt}. Every other reading in this paper holds one
prompt across all seven deployments, including the \pmd{} column beside them,
Table~\ref{tab:hosted-cohort-mode-diversity} and the strict and normalized cell
tables above. Table~\ref{tab:hosted-level-averages-parity} therefore repeats the
comparison over Graph, Countdown, MathIR and Pantry alone, the four domains that share
one prompt. Claude Opus 5 averages
1.42/1.45/1.35 verified modes there at Levels 1--3,
against 1.43/1.44/1.36 over five domains with the
revised Python cell and 1.36/1.17/1.08
with the original one. It is the least varied deployment at every level under all
three readings, and the revised wording raises its Python breadth rather than
lowering it, so the substituted cell moves the five-domain reading away from
concentration rather than towards it. Which of the other six follows it varies
with Python in or out, as their per-domain disagreement in
Table~\ref{tab:hosted-cohort-mode-diversity} already reports; what survives every
reading is the level, not the order.
How far a Python wording change can carry this domain is itself measured, on
three of the other six: App.~\ref{app:sampling-budget-effects}
answers Python under two wordings at 16 prompts and
64 draws, and pooled \pmd{} moves by between
+.024 and +.131 across its
six deployment--level cells, leaving every one of them concentrated
under both wordings (highest \pmd{} .369,
$\neff$ 1.58). That bounds how far wording carries
this domain; it does not identify what this particular revision would do to a
deployment other than the one that answered it.

\begin{table}[H]
  \centering
  \small
  \setlength{\tabcolsep}{5pt}
  \begin{tabular}{lrrrrrr}
    \toprule
    & \multicolumn{2}{c}{Level 1} & \multicolumn{2}{c}{Level 2} & \multicolumn{2}{c}{Level 3} \\
    \cmidrule(lr){2-3}\cmidrule(lr){4-5}\cmidrule(lr){6-7}
    Model & \texttt{pass@8} (\%) & \# modes & \texttt{pass@8} (\%) & \# modes & \texttt{pass@8} (\%) & \# modes \\
    \midrule
    GPT-5.6 Sol & 97.5 & 1.90 & 98.2 & 1.59 & 99.4 & 1.54 \\
    Claude Opus 5 & 97.9 & 1.42 & 99.6 & 1.45 & 99.8 & 1.35 \\
    GPT-5.4 & 97.7 & 2.00 & 98.0 & 1.68 & 99.4 & 1.62 \\
    Grok 4.3 & 99.4 & 2.12 & 100.0 & 1.70 & 100.0 & 1.68 \\
    Kimi K3 & 98.0 & 2.26 & 100.0 & 1.95 & 100.0 & 1.79 \\
    Claude Opus 4.8 & 98.4 & 1.70 & 100.0 & 1.75 & 99.8 & 1.56 \\
    DeepSeek V4 Pro & 98.0 & 2.16 & 98.8 & 2.04 & 99.6 & 1.75 \\
    \bottomrule
  \end{tabular}
  \caption{\textbf{Hosted success and output diversity with Python excluded.}
  The same draws as Table~\ref{tab:hosted-level-averages-medium}, averaged over
  Graph, Countdown, MathIR and Pantry alone: the four domains every deployment answers
  under one prompt, 128 prompts each and eight responses per prompt. Reading this
  table beside that one shows what the revised Python condition of
  App.~\ref{app:hosted-python-prompt} carries, and what it does not.}
  \label{tab:hosted-level-averages-parity}
\end{table}
